Las líneas de continuación más relevantes son:

\begin{itemize}
  \item \textbf{Cierre del Incremento~4}: soporte de polimorfismo
  (\texttt{FunctionalObject} y subclases) mediante \texttt{oneOf} con
  discriminador, y reflejo de los \texttt{relationship()} del ORM como
  \texttt{x-relationships} en el esquema.
  \item \textbf{Codegen}: aprovechar el \emph{bundle} versionado para generar
  automáticamente interfaces TypeScript, validadores y documentación en
  proyectos consumidores.
  \item \textbf{Migración del \texttt{BackendService} heredado} para que
  delegue internamente en \texttt{EntityClient}, retirando progresivamente los
  métodos \emph{hardcoded}.
  \item \textbf{Soporte de \texttt{oneOf/anyOf/allOf}} en el conversor a
  Formly, y \emph{overrides} por \emph{path} (no solo por clave de primer
  nivel).
  \item \textbf{Generación automática de vistas de listado/tabla} con la misma
  filosofía guiada por esquema que ya se aplica a los formularios.
\end{itemize}
